\section{Operational quantities}\label{sec:operational}

Two kinds of numbers are attached to a KS set in the recent literature: the size of the smallest
bipartite Bell scenario with a perfect quantum strategy that it yields, and the graph quantities
of Cabello, Severini and Winter.  Earlier versions reported both with labels (``exact'',
``verified'', ``quantum bound'') that the computations did not justify.  Both have been redone.

\subsection{Bipartite perfect strategies}\label{sec:cabello}

Cabello~\cite{Cabello2024bipartite} relates bipartite perfect quantum strategies to a bipartite
form of KS-uncolorability, and Trandafir and Cabello~\cite{TrandafirCabello2025optimal} convert a
KS set into a strategy with as few measurement settings as possible.  We use the following form.
Let $\mathcal B$ be the set of bases (triads) of a KS set.  Alice holds $S_A\subseteq\mathcal B$,
Bob holds $S_B\subseteq\mathcal B$.  A \emph{local assignment} chooses one vector in each basis
of $S_A$ and one in each basis of $S_B$ such that no vector chosen for Alice is orthogonal to a
vector chosen for Bob.  The pair $(S_A,S_B)$ is \emph{bipartite KS} (B-KS) if no local assignment
exists; the scenario then has $|S_A|$ inputs for Alice and $|S_B|$ for Bob.  A basis held by both
parties is forced to the same choice on both sides, so this agrees with the single-function form
of the definition.

The property is monotone: enlarging $S_A$ or $S_B$ only adds constraints.  The minimum of
$|S_A|\cdot|S_B|$ can therefore be found exactly by the same device as in
Section~\ref{sec:methods}.  A pair that admits a local assignment is enlarged to a maximal such
pair $(M_A,M_B)$, and every B-KS pair must then contain a basis outside $M_A$ on Alice's side or
outside $M_B$ on Bob's.  Sizes $(p,q)$ with $p\le q$ (the definition is symmetric in the
parties) are tried in increasing order of $pq$; the first for which a B-KS pair with
$|S_A|\le p$, $|S_B|\le q$ exists gives the minimum product, every smaller product having been
refuted.  Table~\ref{tab:bpqs} gives the results.  Each run takes at most a few seconds.

\begin{table}[t]
\centering
\caption{Minimum number of input pairs $|S_A|\times|S_B|$ over bipartite-KS pairs of basis
subsets of each set.  Every entry is exact for the set named: all smaller products are excluded.
The first two reproduce the values $5\times9$ and $7\times9$ reported
in~\cite{Cabello2025simplest,TrandafirCabello2025optimal}.  $^{\dagger}$: smallest set found in its pool, not certified to be a minimum.}
\label{tab:bpqs}
\begin{tabular}{lrrcr}
\toprule
Set & Rays & Bases & $|S_A|\times|S_B|$ & Product \\
\midrule
Eisenstein 33 & 33 & 14 & $5\times9$ & 45 \\
Peres 33 & 33 & 16 & $7\times9$ & 63 \\
$\sqrt{-2}$, 33 & 33 & 16 & $7\times9$ & 63 \\
$\CK$ & 31 & 17 & $8\times9$ & 72 \\
$\Q(\sqrt{-7})$, 43 & 43 & 23 & $9\times12$ & 108 \\
Supergolden 47 & 47 & 26 & $10\times16$ & 160 \\
\bpqsextrarows
\bottomrule
\end{tabular}
\end{table}

Three comments.  The values for $\CK$, the $\Q(\sqrt{-7})$ set and the supergolden set were
upper bounds from a heuristic search in version~9; they are now exact.  For $\CK$ this is an
independent reproduction: Trandafir and Cabello already report $8\times9$ as optimal
(\cite{TrandafirCabello2025optimal}, Sec.~IV).  The equality of the
Peres and $\sqrt{-2}$ rows is forced, since the quantity depends only on the orthogonality graph
and the two graphs are isomorphic; it is not an independent observation.  And version~9's
empirical rule $|S_A|+|S_B|\ge|\mathcal B|$ is false: for the $\Q(\sqrt{-7})$ set
$9+12=21<23$, and for the supergolden set $10+16=26$ happens to equal $|\mathcal B|$.

$\CK$ has the fewest rays and is not the cheapest scenario: the Eisenstein set, with two more
rays and three fewer bases, needs 45 input pairs against~72.  That observation is
Cabello's~\cite{Cabello2025simplest}.  The additional computations reported here are the rows below $\CK$; we
have not checked whether any of them has been published.  Nothing
here shows that the alphabet \emph{causes} the cost, and the
statement of version~9 that the islands are ``the algebraic substrates of all perfect quantum
strategies in dimension three'' is withdrawn: Cabello's theorem concerns the bipartite
structure built from a strategy, not coordinate alphabets.

\subsection{Cabello--Severini--Winter quantities}\label{sec:csw}

For a graph $G$ on $n$ vertices, the sum of the probabilities of $n$ events whose exclusivity
graph is $G$ is bounded by the independence number $\alpha(G)$ in noncontextual models, by the
Lov\'asz number $\vartheta(G)$ in quantum theory, and by the fractional packing number
$\alpha^{*}(G)$ under the exclusivity principle alone~\cite{CabelloSeveriniWinter2014}.  The
quantum bound $\vartheta(G)$ is a maximum over \emph{all} sets of vectors with orthogonality
graph containing $G$, in any dimension, and over all states.  It is not the value reached by a
particular set of rays.  For given rays with unit representatives $v_1,\dots,v_n$ in $\C^3$ the largest value of
$\sum_i|\braket{\psi}{v_i}|^2$ over states $\psi$ is the top eigenvalue of the sum of projectors
$\sum_i|v_i\rangle\langle v_i|$, which we call the \emph{attained} value.  Always
$\text{attained}\le\vartheta(G)$.

\begin{table}[t]
\centering
\caption{CSW quantities for the unweighted sum over all rays of each set or pool.  $\alpha$:
exact (MaxSAT).  Attained: top eigenvalue of the sum of the rays' projectors (double precision
from exact coordinates).  $\vartheta$: semidefinite program; the primal and dual optimal values
agree to the digits shown.  $\alpha^{*}$: exact, certified by rational primal and dual solutions
of equal value.  Bold: the rays themselves exceed the noncontextual bound.  $^{\dagger}$: smallest set found in its pool, not certified to be a minimum.}
\label{tab:csw}
\begin{tabular}{lrrrrr}
\toprule
Object & $n$ & $\alpha$ & Attained & $\vartheta$ & $\alpha^{*}$ \\
\midrule
$\CK$ & 31 & 11 & \textbf{11.667} & 11.713 & 12 \\
Peres 33 & 33 & 12 & 11.000 & 12.000 & 12 \\
$\sqrt{-2}$, 33 & 33 & 12 & 11.000 & 12.000 & 12 \\
Eisenstein 33 & 33 & 12 & 11.000 & 12.049 & 13 \\
$\Q(\sqrt{-7})$, 43 & 43 & 16 & 15.000 & 16.000 & 16 \\
Supergolden 47 & 47 & 19 & 15.895 & 19.000 & 19 \\
\cswextrarows
\midrule
pool $\{0,\pm1,\pm2\}$ & 49 & 17 & 16.333 & 17.699 & 19 \\
pool $\{0,\pm1,\pm\sqrt2\}$ & 49 & 23 & 16.333 & 23.000 & 23 \\
pool $\{0,\pm1,\pm\sqrt{-2}\}$ & 49 & 23 & 16.333 & 23.000 & 23 \\
pool $\{0,\pm1,\pm\omega,\pm\omega^2\}$ & 57 & 18 & \textbf{19.000} & 19.340 & 21 \\
pool $\{0,\pm1,\pm\alpha,\pm\alpha^{*}\}$ & 145 & 50 & 48.333 & 55.893 & 66 \\
pool $\{0,\pm1,\pm\psi,\pm\psi^2\}$ & 109 & 41 & 36.333 & 43.952 & 47 \\
\bottomrule
\end{tabular}
\end{table}

Table~\ref{tab:csw} reproduces the values of $\alpha$, $\vartheta$ and $\alpha^{*}$ printed in
version~9 and adds the attained column, which changes the reading of them.  Every raw pool in this table is
invariant under signed permutations of the coordinates, so the sum of its projectors is
$(n/3)$ times the identity and the attained value is $n/3$ for every state.  For $\CK$ the sum is
$\mathrm{diag}(35/3,29/3,29/3)$.

\begin{itemize}
\item Only two of the objects in the table violate the unweighted noncontextuality inequality with
their own rays: $\CK$ ($35/3>11$, for a state along one coordinate axis) and the 57-ray Eisenstein pool
($19>18$, for every state).
\item The $\Q(\sqrt{-7})$ pool has the largest ratio $\vartheta/\alpha=1.118$, which version~9
reported as its ``contextual advantage''.  Its own rays reach $145/3\approx48.3$, below
$\alpha=50$.  The Lov\'asz number of that graph is attained by some other representation; in
what dimension, we have not determined.
\item A KS set violates \emph{some} noncontextuality inequality for every state; the point here
is only that the unweighted sum over all its rays need not be one.
\end{itemize}

Version~9 built on these numbers an account in which ``auxiliary'' rays (those in no triad)
``amplify'' the advantage, supported by a comparison with the Hoffman eigenvalue bound.  Both are
withdrawn.  The Hoffman bound in the form used holds for regular graphs, and these graphs are not
regular; the Peres pool, described there as disconnected, is connected; and the comparison of a
pool with its minimum subset compares two different graphs.  We also avoid the term
``$\vartheta$-perfect'' for a graph with $\vartheta=\alpha$ numerically, since perfection is a
statement about all induced subgraphs.
